%% For double-blind review submission, w/o CCS and ACM Reference (max submission space)
\documentclass[sigconf]{acmart}\settopmatter{printfolios=false,printccs=false,printacmref=false}

%% For double-blind review submission, w/ CCS and ACM Reference
%\documentclass[acmsmall,review,anonymous]{acmart}\settopmatter{printfolios=true}
%% For single-blind review submission, w/o CCS and ACM Reference (max submission space)
%\documentclass[acmsmall,review]{acmart}\settopmatter{printfolios=true,printccs=false,printacmref=false}
%% For single-blind review submission, w/ CCS and ACM Reference
%\documentclass[acmsmall,review]{acmart}\settopmatter{printfolios=true}
%% For final camera-ready submission, w/ required CCS and ACM Reference
%\documentclass[acmsmall]{acmart}\settopmatter{}


%% Conference information
%% Supplied to authors by publisher for camera-ready submission;
%% use defaults for review submission.
\acmConference[]{SIGBOVIK}
\acmYear{}
\acmISBN{} % \acmISBN{978-x-xxxx-xxxx-x/YY/MM}
\acmDOI{} % \acmDOI{10.1145/nnnnnnn.nnnnnnn}
\startPage{1}

%% Copyright information
%% Supplied to authors (based on authors' rights management selection;
%% see authors.acm.org) by publisher for camera-ready submission;
%% use 'none' for review submission.
\setcopyright{none}
%\setcopyright{acmcopyright}
%\setcopyright{acmlicensed}
%\setcopyright{rightsretained}
%\copyrightyear{2018}           %% If different from \acmYear

%% Bibliography style
\bibliographystyle{ACM-Reference-Format}
%% Citation style
%% Note: author/year citations are required for papers published as an
%% issue of PACMPL.
%%\citestyle{acmauthoryear}   %% For author/year citations


%%%%%%%%%%%%%%%%%%%%%%%%%%%%%%%%%%%%%%%%%%%%%%%%%%%%%%%%%%%%%%%%%%%%%%
%% Note: Authors migrating a paper from PACMPL format to traditional
%% SIGPLAN proceedings format must update the '\documentclass' and
%% topmatter commands above; see 'acmart-sigplanproc-template.tex'.
%%%%%%%%%%%%%%%%%%%%%%%%%%%%%%%%%%%%%%%%%%%%%%%%%%%%%%%%%%%%%%%%%%%%%%

%% Some recommended packages.
\usepackage{booktabs}   %% For formal tables:
%% http://ctan.org/pkg/booktabs
\usepackage{multirow}
\usepackage{subcaption} %% For complex figures with subfigures/subcaptions
                        %% http://ctan.org/pkg/subcaption

\usepackage{listings}
%\usepackage{amsmath, amsthm, amssymb, stmaryrd, mathtools, mathrsfs}
\usepackage{mathpartir}
\usepackage{natbib}
\usepackage{url}
%\usepackage{tipa}

\newtheorem{lemma}{Lemma}
\newtheorem{corollary}{Corollary}
\newtheorem{theorem}{Theorem}
\newtheorem{proposition}{Proposition}

\theoremstyle{definition}
\newtheorem{definition}{Definition}
\newtheorem{remark}{Remark}

\newcommand{\jan}[1]{{\color{blue}[[Jan says: #1]]}}
\newcommand{\stefan}[1]{{\color{blue}[[Stefan says: #1]]}}

%\input{macros}

\lstset{
  basicstyle=\ttfamily\footnotesize,
  language=ML
}

\begin{document}

\title{GPTSort: An Earth-Shattering, Paradigm-Shifting New Sorting Algorithm}


\author{Stefan Muller}
\affiliation{\institution{University of Connecticut}}

\begin{abstract}
  We present a new AI-enabled sorting algorithm, GPTSort.
  %
  In contrast to conventional sorting algorithms, GPTSort does not
  require implementers to have any knowledge of integer sorting.
  %
  When run using the GPT 4.1 model, GPTSort can sort~10 integers of up to~7
  digits in under~2 seconds, and is correct almost all of the time.
\end{abstract}

\maketitle

\section{Introduction}
Sorting algorithms, which take a list of comparable items called {\em keys},
and return a permutation of the input list sorted by increasing key,
are some of the first algorithms that undergraduate computer science
students learn.
%
But you don't need to know what they are to appreciate this paper or,
apparently, at all anymore.
%
{\em Algorithms} are predetermined procedural strategies for solving
problems, but apparently you don't need to know that anymore either.
%
You don't need to know these things because Artificial Intelligence (AI)
has now advanced to the point where we can just have it know things for us.
%
So all of these things we used to know are useless because people knew these
things in the past.
%
Things people knew in the past used to be called {\em scholarship}.
%
We now call it {\em training data}.
%
That's on the borderline of things you need to know.
%
But probably not for long.
%
Because soon AI will know that for us too.
%
It used to be that people thought that because AI was the future, it was
important to study how to build and maintain AI tools to make that future
the present, but now that it's the present, those things people knew in the
past ({\em training data}) are now known by the AI and are therefore not
necessary to know anymore.
%
The AI systems of the present were built by people in the past who knew
how to build the AI systems of the present in the future.
%
Nobody in the past knew how to build the AI systems of the present future,
because if they did, those would be the AI systems of the present, but
presumably someone presently knows how to build the AI systems of the future
even though these things aren't important to know anymore.
%
So, wait, who will build the AI systems of the future and how will they know
how to do it?
%
I'm not sure.
%
But it's fine to not be sure about things.
%
Because AI knows them.
%
I should ask it.

What was I talking about?

Right.
%
Sorting algorithms.
%
You don't need to know algorithms for sorting anymore because, like
everything else, AI can do it.
%
Sorting algorithms of the past have included algorithms like
{\em merge sort}, {\em quick sort}, {\em insertion sort}, and {\em bubble sort}.
%
People used to know these things.
%
They are therefore {\em training data}.
%
But AI changed everything, including how we sort.
%
We will never sort the same way again.
%
There's no going back.
%
Or back and then forward and then back again.
%
Only forward.
%
That's kind of the point of sorting, but it's also the point of technological
progress.

Sometimes I wish it weren't.

Anyway, this paper presents GPTSort, the only sorting algorithm you need to
know now, unless someone realizes it's bad, in which case, it will immediately
be replaced by something marginally better that will then be good and the
current one, which will then be the past one, will have always been bad
and the idea will have always been to quickly replace it with something good,
with the past one (currently the current one) now in the past and therefore
just being {\em training data}.
%
Fortunately, this isn't bad, it's good.
%
Take my word for it.
%
Well, actually, take AI's word for it.
%
I used to know things.
%
Now I am just {\em training data}.

\section{Algorithm Description}

In this section, we present the GPTSort algorithm.
%
Fortunately, the algorithm is very simple, and we can present it in very
few words.
%
Because, if all of my messaging clients are any indication, if the algorithm
consisted of more than 12 words, it would need an AI summary.
%
The pseudocode for the algorithm is in Figure~\ref{fig:pseudocode} and
consists of less than~12 words.
%
I know because I asked AI to count them.
%
I used to know how to count, but now that is {\em training data}.
%
My son is two years old and is learning to count from muppets.
%
I keep telling him he doesn't need to know how to count because AI knows that
now.
%
He still likes the muppets.
%
I guess he also doesn't know how to use AI.
%
I guess that's something you still need to know.
%
But doesn't AI know how to use AI, because people used to know how to use
(worse) AI?
%
So why do I still have to know it?
%
My head hurts.
%
I should get a diagnosis from AI for why my head hurts.

\begin{figure}
\begin{lstlisting}[showspaces=false]
Algorithm GPTSort(List L):
  prompt = "Can you sort this list for me? Use a correct
          computational method and output the format as
          a comma-separated list with no other text."
  response = LLM_API.create(prompt + join(',', L))
  return response.split(',')
\end{lstlisting}
\caption{Pseudocode for GPTSort}
\label{fig:pseudocode}
\end{figure}

Anyway, as you can see from Figure~\ref{fig:pseudocode}, the algorithm is
quite simple.
%
First it turns the list into a string prefixed with a short explanation text
called a {\em prompt} explaining what AI knows that I'd like to temporarily
know until I forget and have to ask AI again.
%
Back when we used to learn things from people, we used to call these
{\em questions} and they used to end in these curvy exclamation point things?
%
Just look at those things.
%
Wild.
%
Now they're prompts.
%
Questions are {\em training data}.

We {\em prompt} the AI to sort the list, then we get back the sorted list.
%
Done.

\section{Proof of Correctness}

What's this section doing here?
%
In the algorithms papers that make up {\em training data}, there's a
section called Proof of Correctness where the authors prove that their
algorithm gives the correct answer.
%
That's the training data.
%
Our goal is to make what we know look enough like the training data that
people can't tell them apart.
%
That's how we know we pass the Turing test.
%
Remember when we used to talk about the Turing test?
%
Good times.

So we need this section, and it should say how we know the algorithm is
correct.
%
The reason is obvious: it's right there in the prompt.
%
I told the AI to use a {\em correct} method to sort the numbers.
%
The AI has to do what you say almost all of the time.
%
It's the Third Law of Robotics\cite{irobot}.

\section{Complexity Analysis}
Analyzing the complexity of GPTSort is a little more difficult.
%
It's a constant number of operations, but if you really want to be
pedantic, which is what people used to do, so it's in the {\em training data},
so we have to make it look like we're doing it, one of those operations
(quering a Large Language Model) might do a non-constant amount of work.
%
In order to figure out how much work it's doing, we need to know what it's
doing.
%
I don't really know what an LLM is doing when you ask it something.
%
That's something people used to have to know, but now AI knows it, so I asked
AI.
%
It said it's parsing the input string into a list of numbers by removing the
commas, then using a guaranteed-correct sorting algorithm\footnote{Probably
Timsort, since that's what Python used in the past~\cite{powersort}, and
the past is the training data} and then puts the numbers back together with
commas.
%
Removing and adding commas are presumably done in one pass over the prompt
and the output of the sorting algorithm, so are
both~$O(n)$, and Timsort (according
to AI) is~$O(n\log n)$.

\section{Comparison with Other Sorting Algorithms}\label{sec:compare}
There's not really any need to compare with other algorithms because of
how much of a paradigm shift AI is (it can't possibly be an apples to apples
comparison: it's like comparing apples to way better apples that can think for
themselves and can also postulate the law of gravity themselves rather than
falling on someone who used to know things who could do that himself),
but again, it's in the {\em training data}, so it's expected.
%
See Table~\ref{tab:compare}.

\begin{table*}
  \begin{tabular}{l l l l l l}
    \toprule
    Algorithm & Proven Correct? & Correct? & Average Case & Worst Case & Uses AI\\
    \midrule
    GPTSort & No & No & $n \log n + n\cdot \mathit{poly}(M)$ & Define ``worst'' & Yes\\
      Merge Sort & Yes & Yes & $n \log n$ & $n \log n$ & No\\
      Quick Sort & Yes & Yes & $n \log n$ & $n^2$ & No\\
      Bubble Sort & Yes & Yes & $n^2$ & $n^2$ & No\\
      \bottomrule
  \end{tabular}
  \caption{A comparison of GPTSort with other sorting algorithms.
    $M$ = The size of the LLM (exponential in time, linear in Nvidia's stock price). GPTSort is clearly superior on all of the metrics anyone still cares about.}
  \label{tab:compare}
\end{table*}
    
\section{Experimental Comparison}
I also ran the four sorting algorithms from Section~\ref{sec:compare}.
%
GPTSort was implemented in Python, the other three sorting algorithms were
implemented in OCaml.
%
Really though, the bulk of GPTSort is implemented in some langauge that's
probably much more performant than OCaml, so this comparison should make
GPTSort look pretty good.
%
The GPTSort implementation uses the OpenAI API with GPT 4.1.
%
To make it look even better, we used OCaml's bytecode compiler for the other
three sorting algorithms, and ran them on a several-year-old laptop.
%
We ran each sorting algorithm on lists of 10 and 100 integers between~0
and~$10^7$, with 100 runs for each experiment.
%
Table~\ref{tab:exp-compare} reports the results, with time in microseconds
averaged across all~100 runs, as well as the percentage of the runs in which
each sorting algorithm returned the correct answer.

\begin{table}
  \begin{tabular}{l r r l r r}
    \toprule
    & \multicolumn{2}{c}{$N$=10} & & \multicolumn{2}{c}{$N$=100}  \\
    \cmidrule{2-3} \cmidrule{5-6}
    Algorithm & Time ($\mu$s) & Correct \% & & Time ($\mu$s) & Correct \%\\
    \midrule
    GPTSort & 1,163,740 & 93\% & & 4,669,436 & 3\%\\
    Merge Sort & 4 & 100\% & & 45 & 100\%\\
    Quick Sort & 2 & 100\% & & 40 & 100\%\\
    Bubble Sort & 3 & 100\% & & 372 & 100\%\\
      \bottomrule
  \end{tabular}
  \caption{An empirical comparison of GPTSort with other sorting algorithms.}
  \label{tab:exp-compare}
\end{table}

\paragraph{Discussion of Minor Temporary Deviations from Correctness}
As you may notice in Table~\ref{tab:exp-compare}, the results of GPTSort,
especially for~$N=100$, occasionally differ somewhat from those of
pre-paradigm-shift sorting algorithms.
%
In this subsection, we discuss these deviations.

\textbf{Disclaimer:} Although we have attempted to remain in scope for
SIGBOVIK, this subsection has the potential to contain actual useful
scientific observations.

We ran GPTSort~10 additional times with~$N=100$ and more thoroughly documented
the deviations from correctness.
%
Each run returned incorrect results.
%
Table~\ref{tab:wrong} details, for each of the~10 runs, how many integers
were out of order (i.e., smaller than the previous integer), how many were
duplicated, how many were hallucinated (i.e., appeared in the output list
but not the input list) and how many were dropped (i.e., appeared in the
input list but not the output list).
%
The most common error was placing numbers out of order, with an average
of~3.2 numbers out of order per run.
%
In some cases, numbers were placed in the position they would be in with one
more or fewer digit.
%
This often went along with a duplication or hallucination of an integer.
%
For example, in one run, 795420 appeared twice in the output list, once in
the correct place and once between 7555090 and 8209087.
%
In another run, 272571 appeared in the output list (in the correct position)
but not in the input list, while 2725716 appeared in the input list (and in
the correct position in the output list).
%
In some cases, out-of-order integers are not as easily explainable, but
are usually close to the correct position (for example, in one run,
26202 and 26128 are swapped).
\begin{table}
  \begin{tabular}{l r r r r}
    \toprule
    Run & Out of Order & Duplicated & Hallucinated & Dropped\\
    \midrule
    1 & 7 & 0 & 0 & 2\\
    2 & 4 & 1 & 1 & 0\\
    3 & 3 & 0 & 1 & 2\\
    4 & 1 & 1 & 1 & 0\\
    5 & 7 & 5 & 0 & 0\\
    6 & 1 & 1 & 0 & 0\\
    7 & 6 & 0 & 2 & 0\\
    8 & 2 & 0 & 1 & 0\\
    9 & 1 & 1 & 1 & 1\\
    10 & 0 & 0 & 4 & 0\\
    \midrule
    Total & 32 & 9 & 11 & 5\\
    \bottomrule
  \end{tabular}
  \caption{A detailed look at deviations from correctness in~10 runs
    of GPTSort with~$N = 100$.}
  \label{tab:wrong}
\end{table}

\section{Conclusion}
As you can see, this paper now exists.
%
It is therefore {\em training data} and will be used by AI models to
build better and better sorting algorithms in the future.
%
This can have only positive impacts on scientific discovery and
sorting of integers.

%\bibliographystyle{plainnat}
\bibliography{main}

\end{document}
